\lyseis{Β1.1}{Βασικές έννοιες: σημείο, ευθεία, ημιευθεία, ευθύγραμμο τμήμα, γωνία, κύκλος}
\sol{1} α) σημείου. β) ευθεία. γ) άπειρες, μία και μοναδική. δ) ημιευθεία. ε) αντικείμενες.
\sol{2} Ελεύθερη κατασκευή. Προκύπτουν $3$ ευθείες ($\mathrm{ΑΒ}$, $\mathrm{ΑΓ}$, $\mathrm{ΒΓ}$) και $3$ ευθύγραμμα τμήματα ($\mathrm{ΑΒ}$, $\mathrm{ΑΓ}$, $\mathrm{ΒΓ}$).
\sol{3} α) $\mathrm{ΚΛ}$ \quad β) $\mathrm{Κx}$ (με το $x$ προς το $\Lambda$) \quad γ) η ευθεία $\mathrm{ΚΛ}$.
\sol{4} Ελεύθερη κατασκευή. Κάθε ευθεία χωρίζεται από το $\mathrm{Ο}$ σε δύο αντικείμενες ημιευθείες, άρα βλέπεις $2\cdot2=4$ ημιευθείες συνολικά.
\sol{5} Ελεύθερη κατασκευή. Π.χ. $x\hat{\mathrm{Ο}}y$ και $y\hat{\mathrm{Ο}}x$ (ή $\hat{\mathrm{Ο}}$). Η κυρτή γωνία είναι η μικρότερη περιοχή και η μη κυρτή η μεγαλύτερη.
\sol{6} Ελεύθερη κατασκευή. Και οι τρεις ακτίνες βγαίνουν $2{,}5$ cm: όλες οι ακτίνες του ίδιου κύκλου είναι ίσες.
\sol{7} α) Λάθος: η ευθεία δεν έχει άκρα. β) Σωστό. γ) Λάθος: έχουν διαφορετική αρχή. δ) Σωστό. ε) Σωστό. στ) Λάθος: έχουν το ίδιο κέντρο, αλλά οι ακτίνες τους είναι διαφορετικές.
\sol{8} Αρκούν δύο παλούκια. Τεντώνει το σχοινί ανάμεσά τους, το προεκτείνει προς τις δύο μεριές και φυτεύει κατά μήκος του.
\sol{9} Η δέσμη είναι ημιευθεία με αρχή το $\mathrm{Α}$. Οι δύο φακοί που φωτίζουν ο ένας τον άλλον δίνουν δύο ημιευθείες πάνω στην ίδια ευθεία, με αντίθετες κατευθύνσεις. Αν οι φακοί ήταν στο ίδιο σημείο, θα ήταν αντικείμενες ημιευθείες.
\sol{10} Κάθε ζευγάρι ορίζει ένα σχοινί: $\dfrac{5\cdot4}{2}=10$ σχοινιά.
\sol{11} Τέμνουν το μονοπάτι τα $\mathrm{ΚΜ}$ και $\Lambda\mathrm{Μ}$, γιατί ενώνουν σημεία από διαφορετικά ημιεπίπεδα. Το $\mathrm{ΚΛ}$ δεν το τέμνει.
\sol{12} Ο κύκλος έχει ακτίνα $1{,}2$ μ. Το φυτό στα $1{,}2$ μ. είναι πάνω στη γραμμή του κύκλου. Το φυτό στα $1{,}6$ μ. είναι έξω από τον κύκλο (και αυτό στα $0{,}9$ μ. είναι μέσα).
\sol{13} Η ευθεία είναι άπειρη προς τις δύο μεριές και δεν έχει άκρα, γι' αυτό δεν έχει μήκος. Το τμήμα έχει άκρα και μήκος. Η ημιευθεία είναι άπειρη προς μία μεριά, άρα δεν μετριέται.
\sol{14} Κάθε σημείο ορίζει δύο ημιευθείες, άρα $3\cdot2=6$.
\sol{15} Όχι: από δύο διαφορετικά σημεία περνά μία και μοναδική ευθεία. Αν περνούσαν δύο, δεν θα ήταν διαφορετικές.
\sol{16} Όχι. Ένα κοινό σημείο θα απείχε από το $\mathrm{Ο}$ ταυτόχρονα $3$ cm και $5$ cm, που είναι αδύνατον.
\sol{17} Ελεύθερη απάντηση. Π.χ. σημείο: το ίχνος της μύτης του μολυβιού· τμήμα: η ακμή του τραπεζιού· επίπεδο: το δάπεδο· ημιευθεία: η δέσμη ενός φακού.
\sol{18} Ελεύθερη κατασκευή. Η δεύτερη καρφίτσα ορίζει τη θέση του νήματος και το νήμα δείχνει την ευθεία. Ο κύκλος έχει ως κέντρο το σημείο που διάλεξες και περνά από το άλλο σημείο.
